% SPDX-License-Identifier: MIT
\chapter{Tree file formats}
\label{ch:formats}

\section{What \app{} reads}

\begin{center}
\begin{tabular}{@{}llp{0.48\linewidth}@{}}
\toprule
\textbf{Format} & \textbf{Extensions} & \textbf{Notes} \\
\midrule
Newick   & \cmd{.nwk}, \cmd{.tree}, \cmd{.tre}, \cmd{.newick} &
  Full grammar: single-quoted labels with doubled-quote escapes, nested comments,
  branch lengths, support values. \\
NHX      & \cmd{.nhx} &
  New Hampshire eXtended. \cmd{[\&\&NHX:key=value]} annotations become typed node
  attributes. \\
NEXUS    & \cmd{.nex}, \cmd{.nexus} &
  Including \cmd{TRANSLATE} blocks and \cmd{[\&R]}/\cmd{[\&U]} rooting hints. \\
phyloXML & \cmd{.xml}, \cmd{.phyloxml} &
  Parsed with XML entity expansion disabled. \\
\bottomrule
\end{tabular}
\end{center}

Format is detected from the content, not only the extension, so a Newick tree
saved as \cmd{.txt} still loads. Pass \opt{--format} to the CLI to override
detection.

BEAST and FigTree \cmd{[\&key=value]} metacomments are also parsed, so trees from
those pipelines keep their posterior probabilities and rate annotations.

\section{Writing}

\app{} writes Newick, NEXUS and phyloXML. NHX is read but not written as a
separate format; NHX attributes are preserved on the nodes and can be written
back into Newick comments.

\section{Lenient parsing}

Real-world tree files violate their specifications routinely: unbalanced quotes,
duplicate tip labels, negative branch lengths, support values sitting where an
internal label should be. A parser that refuses such files is a parser you cannot
use on the files you actually have.

\app{} therefore records each recoverable problem as a \emph{diagnostic} and
carries on. Exceptions are reserved for input that cannot be interpreted at all.
Diagnostics are shown to you (\keys{Ctrl+D}) rather than hidden, so leniency does
not become silence.

\subsection{Diagnostics you are likely to see}

\begin{description}[style=nextline,leftmargin=0pt,labelindent=0pt,labelwidth=0pt]
  \item[\cmd{newick.internal-label-ambiguous}]
    Internal labels that parse as numbers. These are almost always bootstrap or
    posterior support values, so they are read as support. If yours are genuinely
    names, pass \cmd{internal\_labels='name'} through the API. This is the single
    most common diagnostic and it is informational.
  \item[\cmd{newick.duplicate-name}]
    Two tips share a name. Both are kept; annotation matching by name will attach
    to both, which is usually not what you want.
  \item[\cmd{newick.negative-length}]
    A negative branch length. By default these are treated as zero for layout
    purposes, which is controlled by \cmd{ignore\_negative\_lengths}.
  \item[\cmd{compose.approximate-band-space}]
    Raised when annotation tracks are drawn in an unrooted layout. See the
    warning in Section~\ref{sec:unrooted-tracks}.
\end{description}

\section{Security}

phyloXML is XML, and XML parsers are a classic vector for two attacks. \app{}
disables entity expansion entirely, which refuses both external-entity injection
(reading files off your disk) and entity-expansion denial of service (the
``billion laughs'' attack). A phyloXML file from an untrusted source cannot use
either against you.

\section{Large files}

Parsing is linear in tree size. A 100{,}000-tip Newick file parses in about
2.2~seconds and lays out in about 0.8~seconds on a normal workstation.

All traversals are iterative rather than recursive, which matters more than it
sounds: a deeply unbalanced tree --- the shape serially sampled sequence data
produces --- would exceed Python's recursion limit at a few thousand nodes under
a recursive implementation. \app{} handles a 50{,}000-node caterpillar tree
without difficulty.

The application remains responsive for trees of a few thousand tips. Beyond that,
switching layout mode takes several seconds during which the window does not
repaint; the work is not currently cancellable.
